% Central caption registry shared by the main paper and Supplementary Information.

\newcommand{\FigureOneCaption}{Mechanism overview on the selected branch.
\textbf{a}, Thin grey contours show the isotropic full-wave local-frequency
surface, and the solid grey curve marks the continuous ZGV ring at the exact
isotropic radius \(\kappa_0\).  \textbf{b}, Angular anomaly of the perturbation: the difference between
the anisotropic and isotropic surfaces with its mean over angle removed at
each radius.  The removed part is the angle-independent carrier shift,
which moves the ring without breaking it, so the fourfold remainder shown
here is the part that breaks it.  Troughs and crests are shaded in the
minimum and saddle colours; the computed, independently refined minima
(blue circles) and saddles (orange diamonds) fall on the troughs and
crests respectively, within the declared local annulus.  Marker shape
duplicates colour, and the panel regroups the two stored surfaces without
introducing a new computation.  \textbf{c}, Conceptual summary of the
change from a one-dimensional critical ring with generic \(t^{-1/2}\) decay to
isolated Morse points with generic \(t^{-1}\) decay; this panel is explanatory
rather than numerical, with the computed asymptotic slices tested in
Fig.~\ref{fig:decay-crossover}.  This mechanism overview does not carry the
numerical realization claim; the finite-resolution full-wave search,
Cartesian Hessians, and stability checks are reported in
Fig.~\ref{fig:morse-points}.  Wavevector components are normalized as
\(k_xh\) and \(k_yh\), frequencies as \(\Omega=\omega h/c_T\), and panels
\textbf{a} and \textbf{b} use common axes.  Source data:
\texttt{data/source\_data/figure\_01}.}

\newcommand{\FigureTwoCaption}{Independent numerical verification of the
isotropic reference.  \textbf{a}, The solid curve is the independently
discretized full-wave symmetric branch; the open orange circular symbol marks
the exact Rayleigh--Lamb reference \((\kappa_0,\Omega_0)\).  \textbf{b}, The
solid curve with circular markers is the local full-wave branch, whereas the
dashed line is the coefficient prediction
\(\Omega_0+a(\kappa-\kappa_0)^2/2\).  The positive curvature \(a\) is obtained
from exact implicit differentiation, not from a fit to the plotted samples.
\textbf{c}, Coloured marker--lines on the left axis report relative
\(\kappa_0\), \(\Omega_0\), and curvature errors together with the normalized
eigen residual, Hermitian defect, and mass-orthogonality defect; the
short-dashed line on the right axis is the relative eigengap.  \textbf{d},
Through-thickness profiles of the mass-normalized mode: solid, short-dashed,
and long-dashed curves denote \(\lvert u_x\rvert\), \(\lvert u_y\rvert\), and
\(\lvert u_z\rvert\), divided by their common global maximum.  The neutral
dotted curve is the squared-displacement proxy \(\sum_i\lvert u_i\rvert^2\),
normalized by its own maximum; it is neither a signed-parity nor a strain-energy
diagnostic.  Here \(\kappa=kh\), \(\Omega=\omega h/c_T\), and \(z/h\) is the
normalized thickness coordinate.  Source data:
\texttt{data/source\_data/figure\_02}.}

\newcommand{\FigureThreeCaption}{Full-elastic sensitivities for the declared
weak cubic family.  \textbf{a}, The solid curve is the analytic full-elastic
coefficient \(V(\theta)-V_0\), and the dashed curve is
\(V_4\cos(4\theta)\).  A common radial offset is used only to render the polar
plot, and the rings are drawn unlabelled; the outermost petal reaches
\(\lvert V_4\rvert\) in units of \(\Omega/\varepsilon\), whose value is
given in panel \textbf{c}.  \textbf{b}, Stems and circular markers show every
stored harmonic amplitude \(\lvert V_m\rvert\) on a logarithmic scale; the
\(m=4\) marker is highlighted and non-fourfold leakage is retained.
\textbf{c}, Filled circles show the stored physical shift
\(\varepsilon V_4\), while the dashed line shows \(Q_4\Delta_C\).  Their
agreement checks the perturbation convention
\(V_4=Q_4\Delta_C^{(1)}\) and
\(\Delta_C=\varepsilon\Delta_C^{(1)}\), rather than providing an independent
full-wave comparison.  \textbf{d}, Solid lines are the analytic full-elastic
\(V\) and \(B=\partial_\kappa V\); open circles and squares are independent
centered full-wave finite differences.  Agreement is quantified by the relative errors
\(\lVert V_{\rm FD}-V\rVert_2/\lVert V\rVert_2\) and
\(\lVert B_{\rm FD}-B\rVert_2/\lVert B\rVert_2\), which the build asserts
to lie below the registered tolerance.  The inset gives the
declared centered-difference step convergence of the relative \(V_4\) and
\(B\) errors.  Source data: \texttt{data/source\_data/figure\_03}.}

\newcommand{\FigureFourCaption}{Full-wave numerical realization of the local
Morse splitting with eight resolved roots.  \textbf{a}, Grey contours show the isotropic full-wave surface,
the solid grey curve marks the exact isotropic critical ring, and the two
short-dashed circles delimit the declared annulus.  \textbf{b}, Filled and
line contours show the anisotropic full-wave surface using the same levels as
panel \textbf{a}; blue circles and orange diamonds are the stored refined
minima and saddles, and the short-dashed circles again show the annular
boundaries.  \textbf{c}, Cartesian Hessian eigenvalues are ordered by
critical-point angle on a symmetric logarithmic axis.  Marker shape and colour
identify Morse class, filled and open symbols denote \(\lambda_1\) and
\(\lambda_2\), vertical bars give the declared Hessian-discrepancy estimate,
and the shaded band spans the maximum such estimate about zero.
\textbf{d}, Filled symbols joined by the solid guide line give the signed
radial error
\((\kappa_{\rm full}-\kappa_{\rm pert})/\kappa_0\); open symbols joined by the
dashed guide line give the signed frequency error
\((\Omega_{\rm full}-\Omega_{\rm pert})/\Omega_0\).  The perturbation values
are first-order predictions, and neither guide line is a fit.  All surfaces
and coordinates use the normalized \(kh\) and \(\omega h/c_T\) convention.
Source data: \texttt{data/source\_data/figure\_04}.}

\newcommand{\FigureFiveCaption}{Perturbation checks over the complete
declared weak-anisotropy sequence.  \textbf{a}, Open circles are the
full-wave minimum--saddle frequency splitting, and the dashed line is its
first-order coefficient prediction.  The reported splitting slope is an
unweighted log--log fit over every declared positive \(\lvert\varepsilon\rvert\),
not a plot-selected range; no separate fitted line is drawn.  \textbf{b}, Blue
circles and orange diamonds show the baseline-corrected full-wave radial shifts
\(q_{\min}/\varepsilon\) and \(q_{\rm sad}/\varepsilon\), while the
corresponding dashed lines are the coefficient limits
\(-B(\theta_j)/a\).  The annotated radial slope is obtained from the
uncompensated quantity \(\max(\lvert q_{\min}\rvert,\lvert q_{\rm sad}\rvert)\)
over the same complete sequence.  \textbf{c}, Markers and connecting guide
lines show the absolute full-wave minus first-order critical-frequency errors
divided by \(\varepsilon^2\) for the minimum and saddle families.  The
reported remainder slope is fitted to the maximum uncompensated error across
the two families; the connecting lines are not fits.  \textbf{d}, Circle and
diamond markers show the minimum and saddle angular orbits for opposite
perturbation signs, demonstrating exchange of the axis and diagonal roles
under \(\varepsilon\rightarrow-\varepsilon\).  Source data:
\texttt{data/source\_data/figure\_05}.}

\newcommand{\FigureSixCaption}{Two numerical asymptotic slices of the
coefficient-resolved response \(G^+\), with \(t=t_{\rm phys}c_T/h\).
\textbf{a}, Full-wave real parts and magnitudes, after removal of the shifted
carrier \(\Omega_0+\varepsilon V_0\) and uniform scaling, collapse onto
\(J_0(\tau)\), \(\tau=\varepsilon V_4t\), for the four declared uniform
transition rows.  The comparison is performed without fitted alignment.
\textbf{b}, Thin curves show the absolute complex error at all stored times;
thick segments show the uniform transition window \(t\geq1500\),
\(\lvert\tau\rvert\leq2\), including Bessel zeros.  \textbf{c}, Pale bands mark the two fitted
intervals.  The early fit
uses the raw \(\varepsilon=0.005\) magnitude, whereas the late fit uses four
complete beat-period RMS values in the fixed-anisotropy
\(\varepsilon=0.08\) row.
These are two slices, not one numerical time trajectory; the proved
growing-\(\lvert\tau\rvert\) overlap connects their coefficients and phases.
\textbf{d}, Crossover markers are first downward crossings of magnitude
\(0.9\) within a theory-centred bracket.  This is a theory-centred consistency
diagnostic; panel \textbf{a} is the primary numerical test.  \textbf{e}, At
\(\varepsilon=0.08\), the full integral and exact-Hessian Morse sum are compared
over \(1500\leq t\leq10200\); open markers denote cancellation samples with
coherence below \(0.30\).  \textbf{f}, Exact minimum and saddle frequencies
mark the critical-frequency features.  Their separation is
\(2\lvert\varepsilon V_4\rvert\) and the modulation magnitude is
\(\lvert\varepsilon V_4\rvert\).  Source data:
\texttt{data/source\_data/figure\_06}.}

\newcommand{\SupplementaryFigureOneCaption}{Polynomial-order and assembly
cross-checks for the selected isotropic branch.  \textbf{a}, Solid
marker--lines give relative errors in the exact Rayleigh--Lamb reference
\(\kappa_0\) (blue circles), \(\Omega_0\) (orange squares), and curvature
\(a\) (dark-blue diamonds) on logarithmic axes.  Each reported fine GLL order
\(p\) pairs a tracked evaluator at \(p-4\) with its independent diagnostic at
\(p\); the machine-precision plateaux need not be monotone.  \textbf{b}, The
same colours and markers show an independently reconstructed one-element
spectrum, while the coloured stars at the final order report the independent
two-element result.  \textbf{c}, At the shared final order, grey circles and
dark-blue stars compare the one- and two-element errors for all three
quantities.  The declared gates are relative \(\kappa_0\) and \(\Omega_0\)
errors below \(10^{-7}\), relative curvature error below \(10^{-4}\), eigen
residual below \(10^{-10}\), and relative eigengap greater than ten times the
larger of that residual and machine precision.  These are local discretization
and mesh-split diagnostics, not a proof for the complete spectrum.  Source
data: \texttt{data/source\_data/supplementary}.}

\newcommand{\SupplementaryFigureTwoCaption}{Quadrature, interpolation, and
phase-discrepancy convergence for the declared response window.  All panels use
the nested angular--radial grids \((N_\theta,N_q)=(16,65),(32,129),(64,257)\)
and a \((128,513)\) reference surface; solid curves with circular markers and
logarithmic ordinates show directly computed diagnostics.  \textbf{a}, Blue
circles give the RMS norm of the coarse-minus-reference demodulated complex
full-wave response divided by the reference-response RMS over the declared
time samples.  \textbf{b}, Blue circles give the maximum absolute frequency
error after periodic-angular and cubic-radial interpolation of each coarse
surface onto the reference grid.  \textbf{c}, Orange circles give the
accumulated inter-resolution phase-discrepancy estimator at the declared
maximum time, including interpolation and coarse/reference nested-order
frequency differences; the black short-dashed line is the declared acceptance
threshold of \(0.05\) rad.  The finest
quadrature error is required to be at most \(0.05\), and successive quadrature
and interpolation ratios at most one half.  These empirical estimators test
consistency across computed resolutions; they do not bound the unknown
continuum error or imply infinite-time phase accuracy.  Source data:
\texttt{data/source\_data/supplementary}.}

\newcommand{\SupplementaryFigureThreeCaption}{Local branch-tracking and
simple-eigenvalue diagnostics.  \textbf{a}, The solid blue curve with circles
is the mass-weighted modal assurance criterion relative to the \(\kappa_0\)
anchor over the declared local wavenumber window; the acceptance threshold
is \(0.99\).  \textbf{b}, The solid orange curve with diamonds is the relative
eigengap of the same tracked branch over that window.  \textbf{c}, On a
logarithmic ordinate, black circles and a solid line report the eigengap and
dark-blue squares and a solid line report the normalized eigen residual over
the polynomial-order sweep.  Acceptance requires residual below \(10^{-10}\)
and gap greater than ten times the larger of residual and machine precision.
Together these computed quantities support branch continuity and separation
only near the selected ZGV point; they neither identify the branch globally
nor certify unrelated modes.  Source data:
\texttt{data/source\_data/supplementary}.}

\newcommand{\SupplementaryFigureFourCaption}{Independent finite-difference
closure of the analytic generalized-eigenvalue sensitivities.  The horizontal
step axes in panels \textbf{a} and \textbf{b} are reversed so refinement runs
from left to right.  \textbf{a}, Blue circles and a solid line show the
centered-difference relative error in \(V_4\) versus
\(\Delta\varepsilon\).  \textbf{b}, Orange circles and a solid line show the
paired-step diagonal error in \(B\), defined as the Euclidean relative error
of the two-component vector at \(\theta=0\) and \(\pi/4\) when
\(\Delta\varepsilon=\Delta q\).  \textbf{c}, The logarithmically coloured
tensor-product heatmap reports the same vector error for every pair
\((\Delta\varepsilon,\Delta q)\); darker cells denote larger error and printed
annotations give the numerical values.  The sensitivity tolerance is \(10^{-4}\):
the final two \(V_4\) steps and the finest two-by-two \(B\) block satisfy it.
These full-wave centered differences are independent of the analytic
sensitivity formula, although they use the same full-wave eigensolver; they
are not fitted data.  Source data:
\texttt{data/source\_data/supplementary}.}

\newcommand{\SupplementaryFigureFiveCaption}{Sensitivity of the computed
response norm to the declared source and spectral window class.  Each ratio
uses the RMS of the demodulated complex full-wave response over the declared
time samples at \(\varepsilon=0.02\), normalized by the baseline
\((R/h,\sigma/k_0)=(0.5,0.15)\).  \textbf{a}, The annotated heatmap covers
three source radii and three super-Gaussian widths; the diverging colour scale
is centered on the normalized baseline value one.  \textbf{b}, Solid curves
with circular markers show the same ratios against \(R/h\), with colour
distinguishing \(\sigma/k_0\); the black horizontal line is the declared
baseline.  The separately computed widest-window boundary weight is required
to remain below the truncation gate \(5\times10^{-3}\).  This sweep supports
response-norm robustness within the displayed source--window class only.  It
does not establish source independence and is not evidence for either the
crossover exponents or the Morse theorem.  Source data:
\texttt{data/source\_data/supplementary}.}

\newcommand{\SupplementaryFigureSixCaption}{Finite-anisotropy silicon stress
test outside the weak-anisotropy theorem, for a [001] plate with
\(C_{11}=165.7\), \(C_{12}=63.9\), and \(C_{44}=79.6\,\mathrm{GPa}\)
(material source \texttt{doi:10.1107/S1600577514004962}).  \textbf{a}, Filled
blue shading and blue contour lines show the full-wave local-frequency surface
normalized as \(\Omega=\omega h/\sqrt{C_{44}/\rho}\); blue circles and orange
diamonds mark four refined minima and four saddles within the declared
annulus.  The shading is not a
separate quantitative colour scale.  \textbf{b}, Points are ordered by stored
index; orange circles and a solid line show the Cartesian Hessian eigenvalue
\(\lambda_1\), blue squares and a solid line show \(\lambda_2\), and the black
line marks zero.  Here marker identity denotes eigenvalue rather than Morse
class: both values are positive at minima and have opposite signs at saddles.
\textbf{c}, On a logarithmic axis, blue circles and a solid line give the
relative tracking gap, while orange diamonds and a solid line give the
stationarity residual \(\lVert\nabla_{\boldsymbol{k}}\Omega\rVert\); this panel
does not display the separately evaluated annular-boundary consistency check.
This is a finite-anisotropy stress test of the numerical search and local
classification only, not evidence for the weak-anisotropy theorem.  Source
data: \texttt{data/source\_data/supplementary}.}
